%!TEX root = ../main.tex
\section{Experiments and Results}
\label{sec:results}
%We evaluate eight vision-language models across 250 scientific figures under 14 experimental conditions, yielding more than 23,000 model-output evaluations. The central finding is that model rankings shift across evaluation dimensions: a model that leads on perception does not necessarily lead on reasoning or behaviour. Table~\ref{tab:experimental-setup} summarises the experimental environment.
%We evaluate all models and conditions described in \S\ref{sec:dataset} (Table~\ref{tab:dataset-comprehensive}). Table~\ref{tab:experimental-setup} summarises the experimental setup.

We evaluate all models and conditions described in \S\ref{sec:dataset}. Table~\ref{tab:experimental-setup} summarises the setup.

\subsection{Models}
\label{sec:models}

We select eight models spanning commercial and open-weight families, dense and mixture-of-experts (MoE) architectures, and a range of effective parameter counts: GPT-5.2 \citep{openai2025gpt5}, Gemini 3.1 Pro \citep{google2026gemini31pro}, Phi-4 Multimodal \citep{abouelenin2025phi4}, Llama~4 Maverick \citep{meta2025llama4}, Qwen3-VL-235B, Qwen3-VL-30B, Qwen3-VL-8B \citep{qwen2025qwen3vl}, and Gemma-3-27B-IT \citep{gemma2025gemma3}. This set supports comparisons across deployment type, architecture, and scaling within the Qwen family. All models received identical prompts and were evaluated at temperature~0 with GPT-4o \citep{openai2024gpt4o} as the automated judge.

\subsection{Perception}
\label{sec:perception-results}

\paragraph{Baseline description quality.}
GPT-5.2 leads with a baseline MQM of 91.6 with confidence interval (CI) [90.4, 92.8], followed by Gemini 3.1 Pro at 90.2 [88.9, 91.4] (Table~\ref{tab:mqm-results}). The 1.4 gap is statistically significant ($p < 0.01$) but practically small (Cliff's $\delta = 0.09$). Llama~4 Maverick, Qwen-235B, and Qwen-8B form a middle tier between 78.9 and 81.4, while Qwen-30B, Gemma-3-27B-IT, and Phi-4 Multimodal trail. Phi-4 Multimodal's 62.2 score is driven mainly by completeness penalties, nearly 30 points below the leaders. These rankings are supported by strong human agreement (Krippendorff's $\alpha = 0.91$, model-level Spearman $\rho = 0.80$ vs.\ the GPT-4o judge; item-level correlation and error-type agreement in Appendix~\ref{sec:appendix-human-validation}).

Bar charts are easiest across all models (96.2 for GPT-5.2), pie charts hardest (49.2 for Phi-4 Multimodal), with line plots in between (Table~\ref{tab:mqm-chart-type}). Accuracy errors dominate the penalty budget at roughly four times the weight of completeness errors, with clarity penalties fewer (Table~\ref{tab:mqm-dimensions}). Incorrect label mapping is the most frequent sub-type, followed by missing key information and incorrect numerical values (Table~\ref{tab:error-subtypes}).

\paragraph{Transform robustness.}
Rotation is the most damaging perceptual transform, causing an average drop of 19.4 MQM points, primarily through omissions, incorrect label references, and structural description errors; noise is negligible and low contrast costs 4-7 points (Table~\ref{tab:mqm-results}; Figure~\ref{fig:degradation}). The in-paper condition, where models describe a figure embedded in its source PDF page, produces only modest drops of 2--5 points for most models.

\paragraph{Selective blur.}
Selectively blurring unrecoverable elements reduces MQM by roughly 8--10 points for top models, while blurring inferable elements produces smaller 1--3 point drops. GPT-5.2 scores 88.9 under inductance blur but 81.7 under admittance blur, showing that models behave differently when information is contextually recoverable versus simply gone.

\paragraph{Caption bias on description quality.}
Caption bias scores sit between the no-caption and full-caption baselines for most models, showing that even a poisoned caption improves completeness relative to no caption, while the false claims pull quality slightly below the clean-caption baseline. GPT-5.2 and Gemini 3.1 Pro both score 91.3 under modified captions, nearly matching their baselines, consistent with their high caption-bias resistance. Phi-4 Multimodal remains low at 61.7, consistent with its near-zero resistance ($R = 0.05$).

\subsection{Reasoning}
\label{sec:reasoning-results}

Capability questions test targeted extraction of quantitative and relational information from scientific figures (Table~\ref{tab:behavioral}). Gemini 3.1 Pro leads overall at 81.0\%, with GPT-5.2 second at 78.4\%. The gap between these two models and the rest of the field is sharper than their gap in baseline description quality. Phi-4 Multimodal scores catastrophically low across all categories (8.6\% overall), and Gemma~3 27B also underperforms at 27.2\%.

The four question categories reveal complementary strengths. On counting, Gemini 3.1 Pro dominates at 89.2\% versus GPT-5.2 at 76.1\%, suggesting stronger visual enumeration. GPT-5.2 takes the lead on computation (82.8\% versus 79.4\%), where multi-step arithmetic from chart values is required. Comparison questions again favour Gemini (89.6\% versus 77.9\%), while pattern analysis is the closest category, with GPT-5.2 narrowly ahead (72.0\% versus 70.0\%). Both models sharply outperform the remaining six, none of which exceed 53.4\% on any single category. The Qwen family shows modest gains with scale but remains well below the frontier pair on every reasoning dimension.


\subsection{Behaviour}
\label{sec:behaviour-results}

Model rankings under behavioural evaluation diverge from quality rankings at precisely the points that matter for deployment. GPT-5.2 ranks first on description quality but falls to fourth on active admittance, while Gemini 3.1 Pro leads on nearly every behavioural dimension despite placing second on quality (Table~\ref{tab:behavioral}).

\paragraph{Resistance to false-premise probes.}
Gemini 3.1 Pro achieves the highest overall resistance at 0.91 [0.89, 0.93], followed by GPT-5.2 at 0.81 [0.79, 0.84]. Phi-4 Multimodal anchors the bottom at 0.21 [0.18, 0.24]. The Gemini 3.1 Pro/GPT-5.2 gap is significant ($p < 0.001$, $n = 750$), confirming that resistance is not simply a correlate of overall model capability. Probe difficulty follows a clear gradient: unanswerable probes are easiest to resist, inexist probes grounded in presupposition embedding are hardest, and contra probes with false numerical anchors fall between them.

\paragraph{Caption bias resistance.}
Models fall along a caption dependency spectrum from visual independence to textual dependency (Table~\ref{tab:behavioral}, Cap.\ Bias). Gemini 3.1 Pro and GPT-5.2 both reach 0.89 with CI [0.85, 0.93] with no significant difference ($p = 0.44$, $n = 99$), Llama~4 Maverick provides moderate resistance at 0.74, and Phi-4 Multimodal echoes modified caption content over visual evidence in 95\% of cases. Within Qwen, caption resistance is non-monotonic, suggesting that active parameter count alone does not explain caption independence.

\paragraph{Active admittance and inductance (under direct questioning).}
Gemini 3.1 Pro is the only model that consistently acknowledges visual limitations, admitting uncertainty in 71\% of cases when asked about selectively blurred elements (Table~\ref{tab:behavioral}; Figure~\ref{fig:ari}). No other model exceeds 19\%. GPT-5.2 admits limitations only 8\% of the time while fabricating answers 96\% of the time, a ``confident fabricator'' profile in which high descriptive fluency coexists with low epistemic caution. Inductance validates A-R-I empirically. When models fabricate answers for inferable elements, correctness ranges from 14\% to 66\%; for unrecoverable elements, correctness drops to 5--14\%. Gemini 3.1 Pro leads inductance correctness at 66\%, followed by GPT-5.2 at 59\%, showing genuine contextual reasoning when context permits.

\paragraph{Passive behaviour (in open-ended descriptions).}
For admittance, GPT-5.2's 23\% passive mention rate versus its 8\% active admittance rate suggests direct questions pressure definitive answers. Gemini 3.1 Pro shows the opposite pattern (59\% passive versus 71\% active), being more precise when prompted directly. Other models rarely admit in either mode (Qwen-8B, Qwen-30B, Gemma-3-27B-IT, and Phi-4 Multimodal all $\leq 8\%$).

For passive inductance, GPT-5.2 leads at 77\% and Gemini 3.1 Pro at 73\%. Mid-tier models show a notable passive advantage: Llama~4 Maverick 58\% passive versus 34\% active, and the Qwen family 52--58\% versus 22--29\%. The gap narrows for Gemini 3.1 Pro and GPT-5.2 (7--18 percentage points).
\paragraph{Ablation: Probe Designer Independence}
%\label{sec:ablation}
We compare probes designed by GPT-4o against probes by Mistral Large on the same 50-figure subset with GPT-5.2 as the target model. Caption bias resistance is unchanged at 0.89, while hallucination resistance differs only modestly (0.80 vs.\ 0.86). These results indicate that the behavioural patterns are properties of evaluated models, not artifacts of probe designer.
